% =============================================
% File: chapters/theoretical_background.tex
% Chương 1: Cơ sở lý thuyết
% =============================================

\chapter{CƠ SỞ LÝ THUYẾT}
\label{chap:cosolythuyet}

Chương này trình bày các kiến thức nền tảng phục vụ cho quá trình nghiên cứu,
bao gồm: đặc tính của động cơ DC NF5475, nền tảng vi điều khiển ESP32, lý thuyết
bộ điều khiển PID, mạng nơron nhân tạo và động học cơ cấu 5-bar parallel.


% ==============================================
\section{Động cơ DC NF5475}
% ==============================================

\subsection{Tổng quan}

Động cơ NF5475 (NISCA NF5475F) là động cơ DC Brushed có tích hợp encoder quang học
phía sau trục động cơ. Đây là loại động cơ được tháo từ máy photocopy thương hiệu
Cannon, Epson, Brother, Ricoh, nhờ đó có chất lượng cao với mức giá phù hợp.

Các ưu điểm nổi bật của động cơ NF5475:
\begin{itemize}
    \item Độ ổn định cao, tiếng ồn nhỏ.
    \item Độ nhiễu gây ra bởi động cơ lên mạch điều khiển thấp.
    \item Chất lượng xung Encoder A, B hoàn hảo.
    \item Chống phát nóng nhiều trong quá trình sử dụng.
\end{itemize}

\subsection{Thông số kỹ thuật}

Các thông số điện cơ chính của NF5475 được tóm tắt trong Bảng~\ref{tab:nf5475_spec}.

\begin{table}[H]
\centering
\caption{Thông số kỹ thuật động cơ NF5475}
\label{tab:nf5475_spec}
\renewcommand{\arraystretch}{1.3}
\begin{tabular}{|l|l|}
\hline
\textbf{Thông số} & \textbf{Giá trị} \\
\hline
Điện áp hoạt động & 12 -- 24 V \\
Công suất định mức & 33 W \\
Dòng tải định mức & 1,49 A \\
Dòng không tải & 0,27 A \\
Tốc độ không tải & 3818 RPM \\
Momen khởi động & 464 mN·m \\
Đường kính trục & 8 mm \\
Đường kính thân động cơ & 54 mm \\
Chiều dài (kể cả encoder) & 93 mm \\
\hline
\end{tabular}
\end{table}

\subsection{Encoder}

Encoder quang học tích hợp trên NF5475 có các đặc điểm:
\begin{itemize}
    \item Điện áp cấp: 5 VDC, 4 dây kết nối.
    \item Số kênh: 2 kênh A, B (A, B pulse channel).
    \item Độ phân giải: 200 xung/vòng (PPR).
    \item Tốc độ phản hồi xung: 20 kHz.
    \item Loại: Optical, Incremental.
\end{itemize}

Nguyên lý xác định chiều quay: khi quay thuận, kênh A sớm pha hơn kênh B;
khi quay nghịch, kênh B sớm pha hơn kênh A.


% ==============================================
\section{Vi điều khiển ESP32 WROOM}
% ==============================================

\subsection{Giới thiệu}

ESP32 WROOM-32 là module vi điều khiển 32-bit của Espressif Systems, tích hợp
Wi-Fi và Bluetooth, được lập trình qua Arduino IDE. Module có đặc điểm phù hợp
với ứng dụng điều khiển nhúng:

\begin{itemize}
    \item CPU: Xtensa LX6 dual-core, 240 MHz.
    \item Flash: 4 MB, RAM: 520 KB.
    \item Hỗ trợ PWM phần cứng (LEDC) với độ phân giải 8-bit, tần số 20 kHz.
    \item Hỗ trợ ngắt ngoài (GPIO interrupt) để đọc encoder.
\end{itemize}

\subsection{Cấu hình chân sử dụng}

Trong đề tài này, ESP32 được cấu hình sử dụng các chân sau:
\begin{itemize}
    \item GPIO 18, 19: chân PWM điều khiển chiều quay động cơ (R\_PWM, L\_PWM).
    \item GPIO 25, 26: chân tín hiệu encoder (phase\_a, phase\_b).
    \item 5V0: cấp nguồn 5V cho encoder.
    \item GND: nối đất chung.
\end{itemize}

Chân \texttt{phase\_a} được cấu hình là chân interrupt chế độ 1x
(chỉ xét trên một cạnh của xung A). PWM xuất ra 8-bit với dải $[-255, 255]$.


% ==============================================
\section{Bộ điều khiển PID}
% ==============================================

\subsection{Nguyên lý hoạt động}

Bộ điều khiển PID (Proportional--Integral--Derivative) là bộ điều khiển vòng kín
có phản hồi, điều khiển ngõ ra bám theo tín hiệu đặt (setpoint). Công thức tổng quát:

\begin{equation}
    u(t) = K_P e(t) + K_I \int_0^t e(\tau)\,d\tau + K_D \frac{d}{dt}e(t)
    \label{eq:pid}
\end{equation}

trong đó $e(t) = r(t) - y(t)$ là sai số giữa setpoint và ngõ ra thực tế.

\subsection{Vai trò từng khâu}

\begin{itemize}
    \item \textbf{Khâu P ($K_P$):} Tỉ lệ với sai số hiện tại. $K_P$ càng lớn,
    đáp ứng càng nhanh nhưng độ vọt lố tăng và có thể mất ổn định.

    \item \textbf{Khâu I ($K_I$):} Tích phân sai số theo thời gian, loại bỏ sai
    số xác lập. $K_I$ càng lớn, sai số xác lập càng nhỏ nhưng độ vọt lố tăng
    và đáp ứng quá độ chậm hơn.

    \item \textbf{Khâu D ($K_D$):} Đạo hàm sai số, dự đoán xu hướng thay đổi.
    $K_D$ giúp giảm vọt lố và ổn định hệ nhanh hơn, nhưng nhạy cảm với nhiễu.
\end{itemize}

\subsection{Mô hình điều khiển động cơ trong đề tài}

Trong đề tài, cấu trúc điều khiển sử dụng hai vòng lặp:

\begin{itemize}
    \item \textbf{Vòng trong:} Bộ PI vận tốc — bám vận tốc đặt ($K_D = 0$).
    \item \textbf{Vòng ngoài:} Bộ PID vị trí — sinh ra vận tốc đặt cho vòng trong.
\end{itemize}

Bộ PI vận tốc được cài đặt với anti-windup khâu I (giới hạn tích phân trong
$[-1000, 1000]$) và giới hạn PWM trong $[-255, 255]$ để bảo vệ động cơ.


% ==============================================
\section{Mạng Nơron Nhân Tạo (ANN)}
% ==============================================

\subsection{Tổng quan}

Mạng Nơron Nhân Tạo (Artificial Neural Network -- ANN) là mô hình tính toán
lấy cảm hứng từ hệ thần kinh sinh học, có khả năng học ánh xạ phi tuyến từ
đầu vào đến đầu ra thông qua dữ liệu huấn luyện. Trong đề tài, ANN được dùng
cho bài toán \textbf{hồi quy} (regression): dự đoán giá trị $K_p$, $K_i$ từ
trạng thái sai số.

\subsection{Hàm kích hoạt}

\begin{itemize}
    \item \textbf{ReLU (Rectified Linear Unit):} Giá trị âm được thay bằng 0,
    giá trị dương giữ nguyên. Phù hợp cho các lớp ẩn trong bài toán hồi quy,
    tránh vấn đề vanishing gradient.
    \begin{equation}
        f(x) = \max(0, x)
    \end{equation}

    \item \textbf{Linear:} Giá trị đi qua không thay đổi. Dùng ở lớp đầu ra
    để cho phép mạng dự đoán giá trị liên tục không bị giới hạn.
    \begin{equation}
        f(x) = x
    \end{equation}
\end{itemize}

\subsection{Lớp Dense với Regularization}

Mỗi lớp Dense thực hiện phép biến đổi tuyến tính:
\begin{equation}
    \mathbf{y} = \mathbf{W}\mathbf{x} + \mathbf{b}
\end{equation}

Để chống overfitting, lớp Dense được trang bị L2 regularization
(phạt trọng số lớn trong hàm loss) và Dropout (ngẫu nhiên loại bỏ một số
nơron trong mỗi bước huấn luyện).

\subsection{Hàm loss và Optimizer}

\begin{itemize}
    \item \textbf{Hàm loss MSE (Mean Squared Error):}
    \begin{equation}
        \mathcal{L} = \frac{1}{n}\sum_{i=1}^{n}(y_i - \hat{y}_i)^2
    \end{equation}
    Đây là hàm loss chuẩn cho bài toán hồi quy.

    \item \textbf{Optimizer Adam:} Kết hợp Momentum và RMSProp, tự động hiệu
    chỉnh learning rate cho từng tham số. Trong đề tài sử dụng:
    learning rate $= 2 \times 10^{-3}$, decay $= 3 \times 10^{-4}$.
\end{itemize}


% ==============================================
\section{Cơ cấu 5-bar Parallel}
% ==============================================

\subsection{Giới thiệu}

Cơ cấu 5-bar parallel là cơ cấu phẳng khép kín gồm 5 thanh nối, trong đó 2
khớp chủ động ($A_1$, $A_2$) được truyền động bởi 2 động cơ độc lập.
So với cơ cấu nối tiếp 3 bậc tự do, cơ cấu 5-bar parallel giúp:

\begin{itemize}
    \item Giảm số động cơ cần điều khiển (từ 3 xuống 2 mỗi chân).
    \item Tăng độ cứng vững cơ học.
    \item Giảm momen quán tính, nâng cao khả năng phản hồi.
\end{itemize}

\subsection{Thông số hình học}

\begin{table}[H]
\centering
\caption{Thông số hình học cơ cấu 5-bar parallel}
\label{tab:5bar_params}
\renewcommand{\arraystretch}{1.3}
\begin{tabular}{|l|c|c|}
\hline
\textbf{Thông số} & \textbf{Ký hiệu} & \textbf{Giá trị} \\
\hline
Khoảng cách giữa hai khớp chủ động & $A_1A_2$ & 105 mm \\
Chiều dài thanh chủ động & $L_{act}$ & 100 mm \\
Chiều dài thanh bị động  & $L_{pas}$ & 180 mm \\
\hline
\end{tabular}
\end{table}

\subsection{Động học thuận (Forward Kinematics)}

Chọn hệ tọa độ phẳng với gốc tại trung điểm $A_1A_2$:
$A_1 = (-d/2,\,0)$, $A_2 = (d/2,\,0)$ với $d = 105$ mm.

Từ hai góc chủ động $\theta_1$, $\theta_2$, vị trí hai khớp bị động:
\begin{align}
    B_1 &= A_1 + L_{act}[\cos\theta_1,\;\sin\theta_1]^T \\
    B_2 &= A_2 + L_{act}[\cos\theta_2,\;\sin\theta_2]^T
\end{align}

Điểm cuối $P(x,y)$ là giao điểm của hai đường tròn tâm $B_1$, $B_2$,
bán kính $L_{pas}$. Nghiệm được chọn theo nhánh làm việc thực tế của chân robot.

\subsection{Động học nghịch (Inverse Kinematics)}

Với vị trí mong muốn $P_d(x_d, y_d)$, góc chủ động tính từ bài toán tam giác:
\begin{align}
    \theta_1 &= \text{atan2}(y_d,\;x_d + d/2) \pm \arccos\!\left(\frac{L_{act}^2 + r_1^2 - L_{pas}^2}{2 L_{act} r_1}\right) \\
    \theta_2 &= \text{atan2}(y_d,\;x_d - d/2) \mp \arccos\!\left(\frac{L_{act}^2 + r_2^2 - L_{pas}^2}{2 L_{act} r_2}\right)
\end{align}

trong đó:
\begin{align*}
    r_1 &= \sqrt{(x_d + d/2)^2 + y_d^2} \\
    r_2 &= \sqrt{(x_d - d/2)^2 + y_d^2}
\end{align*}

Dấu $\pm$ được chọn theo cấu hình lắp thực tế để tránh nhảy nghiệm và
tránh vùng kỳ dị.

\subsection{Ma trận Jacobian}

Jacobian mô tả quan hệ giữa vận tốc khớp và vận tốc điểm cuối:
\begin{equation}
    \begin{bmatrix} \dot{x} \\ \dot{y} \end{bmatrix}
    = \mathbf{J}(\theta_1, \theta_2)
    \begin{bmatrix} \dot{\theta}_1 \\ \dot{\theta}_2 \end{bmatrix}
\end{equation}

Ma trận Jacobian được dùng để:
\begin{itemize}
    \item Sinh vận tốc đặt cho từng động cơ khi robot cần dịch chuyển điểm
    cuối $P$ theo quỹ đạo mong muốn.
    \item Xác định các vùng kỳ dị (singular configuration), nơi sai số nhỏ
    tại điểm cuối có thể gây vận tốc hoặc momen khớp rất lớn.
\end{itemize}


% ==============================================
\section{Học tăng cường (Reinforcement Learning)}
% ==============================================

\subsection{Giới thiệu}

Học tăng cường (Reinforcement Learning -- RL) là kỹ thuật học máy huấn luyện
phần mềm đưa ra quyết định tối ưu thông qua cơ chế thử-và-sai kết hợp hệ thống
khen thưởng -- trừng phạt. Tương tự cách con người học tập, thuật toán RL tự
khám phá chiến lược tốt nhất, thậm chí chấp nhận hy sinh lợi ích ngắn hạn
(trì hoãn khen thưởng) để đạt mục tiêu dài hạn.

Hệ thống RL hoạt động dựa trên sự tương tác khép kín giữa hai thành phần:
\begin{itemize}
    \item \textbf{Agent (tác nhân):} Đối tượng học tập và ra quyết định, quan
    sát môi trường và chọn hành động nhằm tối đa hóa phần thưởng tích lũy.
    \item \textbf{Environment (môi trường):} Thế giới mà agent tồn tại và tương
    tác, phản hồi lại hành động bằng trạng thái mới và phần thưởng (Reward).
\end{itemize}

Tại mỗi bước thời gian, quy trình diễn ra theo thứ tự:
\begin{enumerate}
    \item Agent quan sát trạng thái hiện tại của môi trường.
    \item Agent chọn hành động dựa trên chính sách hiện tại.
    \item Môi trường chuyển sang trạng thái mới và trả về phần thưởng.
    \item Agent cập nhật chính sách dựa trên tín hiệu phần thưởng nhận được.
\end{enumerate}

\subsection{Các thành phần cốt lõi}

\textbf{Không gian trạng thái và quan sát.}
State $s$ là bức tranh đầy đủ về thế giới tại một thời điểm.
Observation $o$ là những gì agent thực sự nhìn thấy (có thể không đầy đủ).

\textbf{Không gian hành động (Action Space).}
Tập hợp tất cả hành động hợp lệ trong môi trường, gồm hai loại:
\begin{itemize}
    \item \textit{Discrete Action Space:} số hành động hữu hạn.
    \item \textit{Continuous Action Space:} số hành động vô hạn (liên tục).
\end{itemize}

\textbf{Chính sách (Policy $\pi$).}
Policy là quy tắc để agent quyết định hành động tiếp theo dựa trên quan sát:
\begin{itemize}
    \item \textit{Deterministic Policy:} Luôn chọn đúng một hành động cho mỗi
    trạng thái.
    \item \textit{Stochastic Policy:} Chọn hành động theo phân phối xác suất.
    Ký hiệu:
    \begin{equation}
        a_t \sim \pi_\theta(\cdot \mid s_t)
        \label{eq:stochastic_policy}
    \end{equation}
    trong đó $\theta$ là tập tham số của policy.
\end{itemize}

\textbf{Quỹ đạo (Trajectory).}
Chuỗi các sự kiện trong một episode:
\begin{equation}
    \tau = (s_0, a_0, s_1, a_1, \ldots)
\end{equation}
Trạng thái ban đầu lấy mẫu từ phân phối khởi tạo $s_0 \sim \rho_0(\cdot)$.

\textbf{Phần thưởng và Tổng lợi nhuận.}
Phần thưởng tức thời tại bước $t$:
\begin{equation}
    r_t = R(s_t, a_t, s_{t+1})
\end{equation}

Tổng lợi nhuận có chiết khấu (Infinite-horizon discounted return):
\begin{equation}
    R(\tau) = \sum_{t=0}^{\infty} \gamma^t r_t, \quad \gamma \in [0,1)
\end{equation}
Hệ số chiết khấu $\gamma$ đảm bảo tổng hội tụ và ưu tiên phần thưởng gần hơn.

\textbf{Hàm giá trị (Value Functions).}
\begin{align}
    V^\pi(s) &= \mathbb{E}_{\tau \sim \pi}\bigl[R(\tau) \mid s_0 = s\bigr]
    \label{eq:Vpi}\\
    Q^\pi(s,a) &= \mathbb{E}_{\tau \sim \pi}\bigl[R(\tau) \mid s_0=s,\, a_0=a\bigr]
    \label{eq:Qpi}\\
    V^*(s) &= \max_\pi\, V^\pi(s) \label{eq:Vstar}\\
    Q^*(s,a) &= \max_\pi\, Q^\pi(s,a) \label{eq:Qstar}
\end{align}

Quan hệ giữa $V^\pi$ và $Q^\pi$:
\begin{align}
    V^\pi(s) &= \mathbb{E}_{a\sim\pi}\bigl[Q^\pi(s,a)\bigr] \\
    V^*(s)   &= \max_a Q^*(s,a)
\end{align}

\textbf{Phương trình Bellman.}
\begin{equation}
    V^\pi(s) = \mathbb{E}_{a\sim\pi,\, s'\sim P}
    \bigl[r(s,a) + \gamma\, V^\pi(s')\bigr]
    \label{eq:bellman}
\end{equation}
Đây là cơ sở để cập nhật mạng Critic trong kiến trúc Actor-Critic của PPO.

\textbf{Hàm lợi thế (Advantage Function).}
\begin{equation}
    A^\pi(s,a) = Q^\pi(s,a) - V^\pi(s)
    \label{eq:advantage}
\end{equation}
$A > 0$ nghĩa là hành động tốt hơn trung bình (tăng xác suất chọn);
$A < 0$ nghĩa là hành động tệ hơn (giảm xác suất chọn).

\textbf{Mục tiêu tối ưu.}
\begin{equation}
    \pi^* = \arg\max_\pi\, J(\pi),
    \quad J(\pi) = \mathbb{E}_{\tau \sim \pi}\bigl[R(\tau)\bigr]
    \label{eq:rl_objective}
\end{equation}


% ==============================================
\section{Thuật toán PPO}
% ==============================================

\subsection{Khái niệm}

Proximal Policy Optimization (PPO), do OpenAI giới thiệu năm 2017, là thuật
toán on-policy thuộc nhóm Policy Gradient. PPO giải quyết vấn đề bước cập nhật
quá lớn bằng cách giới hạn biên độ thay đổi giữa chính sách mới và cũ, đảm bảo
quá trình học ổn định và tránh sụp đổ hiệu suất.

\subsection{Kiến trúc Actor-Critic}

PPO sử dụng kiến trúc Actor-Critic:
\begin{itemize}
    \item \textbf{Actor ($\pi_\theta$):} Mạng nơron đưa ra hành động.
    \item \textbf{Critic ($V_\phi$):} Mạng nơron đánh giá giá trị trạng thái,
    dự đoán $V^\pi(s)$.
\end{itemize}

\subsection{Tỷ lệ xác suất và GAE}

\textbf{Tỷ lệ xác suất (Probability Ratio):}
\begin{equation}
    r_t(\theta) = \frac{\pi_\theta(a_t \mid s_t)}{\pi_{\theta_\text{old}}(a_t \mid s_t)}
    \label{eq:prob_ratio}
\end{equation}
Khi $r_t(\theta) > 1$: policy mới ưa thích hành động đó hơn policy cũ;
khi $r_t(\theta) < 1$: ít ưa thích hơn; khi $r_t(\theta) = 1$: không thay đổi.

\textbf{Generalized Advantage Estimation (GAE).}
Sai số TD (Temporal Difference error):
\begin{equation}
    \delta_t = r_t + \gamma\, V(s_{t+1}) - V(s_t)
    \label{eq:td_error}
\end{equation}

Ước lượng lợi thế:
\begin{equation}
    \hat{A}_t = \sum_{k=0}^{\infty} (\gamma\lambda)^k\, \delta_{t+k}
    \label{eq:gae}
\end{equation}
Thường chọn $\lambda \approx 0{,}95$ để cân bằng phương sai -- độ chệch.

\subsection{Hàm mục tiêu Clipped Surrogate}

\begin{equation}
    L^{\text{CLIP}}(\theta) = \hat{\mathbb{E}}_t\!\left[
        \min\!\left(
            r_t(\theta)\,\hat{A}_t,\;
            \text{clip}\!\left(r_t(\theta),\,1{-}\varepsilon,\,1{+}\varepsilon\right)\hat{A}_t
        \right)
    \right]
    \label{eq:ppo_clip}
\end{equation}

trong đó $\varepsilon \in [0{,}1;\; 0{,}2]$ là hệ số clipping.

Cơ chế clipping:
\begin{itemize}
    \item Nếu $\hat{A}_t > 0$ (hành động tốt): tăng xác suất nhưng không vượt
    quá $1 + \varepsilon$ để tránh tự tin quá mức.
    \item Nếu $\hat{A}_t < 0$ (hành động tệ): giảm xác suất nhưng không nhỏ hơn
    $1 - \varepsilon$ để duy trì khả năng khám phá.
\end{itemize}

\subsection{Hàm mất mát tổng quát}

\begin{equation}
    L^{\text{PPO}}(\theta) = \hat{\mathbb{E}}_t\!\left[
        L^{\text{CLIP}}(\theta)
        - c_1\, L^{\text{VF}}(\theta)
        + c_2\, S\!\left[\pi_\theta\right](s_t)
    \right]
    \label{eq:ppo_total_loss}
\end{equation}

trong đó:
\begin{itemize}
    \item $L^{\text{VF}}(\theta) = \bigl(V_\theta(s_t) - V_t^{\text{target}}\bigr)^2$:
    hàm mất mát MSE của mạng Critic, hệ số $c_1 = 0{,}5$.
    \item $S[\pi_\theta](s_t)$: entropy bonus khuyến khích agent khám phá,
    hệ số $c_2 \approx 0{,}01$.
\end{itemize}


% ==============================================
\section{Nền tảng mô phỏng}
% ==============================================

\subsection{NVIDIA Isaac Sim}

NVIDIA Isaac Sim là nền tảng mô phỏng robot được phát triển bởi NVIDIA trên
nền tảng Omniverse, dùng để thiết kế, mô phỏng, kiểm thử và đánh giá hệ thống
robot trong môi trường ảo có độ chân thực cao.

Các đặc điểm nổi bật:
\begin{itemize}
    \item Mô phỏng vật lý độ chính xác cao: động lực học, va chạm, ma sát,
    chuyển động cơ cấu.
    \item Hỗ trợ mô phỏng cảm biến: camera, LiDAR, IMU, cảm biến chiều sâu.
    \item Nhập mô hình robot từ CAD, URDF, MJCF và các định dạng 3D khác.
    \item Kết nối ROS/ROS2 qua bridge extension.
    \item Tạo dữ liệu tổng hợp (synthetic data) cho bài toán thị giác máy tính.
\end{itemize}

Trong đề tài, Isaac Sim được sử dụng để:
\begin{enumerate}
    \item Kiểm tra mô hình URDF của robot (vị trí khâu, hướng trục khớp,
    thông số va chạm, hệ số ma sát bánh xe).
    \item Thiết lập môi trường mô phỏng trước khi đưa sang Isaac Lab huấn luyện.
\end{enumerate}

\subsection{NVIDIA Isaac Lab}

NVIDIA Isaac Lab là framework mã nguồn mở xây dựng trên Isaac Sim, chuyên phục
vụ nghiên cứu và huấn luyện các thuật toán AI cho robot, đặc biệt là học tăng
cường (RL).

Các đặc điểm chính:
\begin{itemize}
    \item Mô phỏng song song nhiều môi trường trên GPU, tăng tốc huấn luyện
    đáng kể.
    \item Hỗ trợ nhiều loại robot: di động, bốn chân, hình người, tay máy.
    \item Tích hợp với các thư viện RL: RSL-RL, RL-Games, Stable-Baselines.
    \item Kế thừa toàn bộ khả năng mô phỏng vật lý và cảm biến của Isaac Sim.
\end{itemize}


% ==============================================
\section{Quy trình chuẩn bị mô hình huấn luyện}
% ==============================================

Quá trình đưa mô hình robot vào môi trường huấn luyện RL trải qua các bước:

\begin{enumerate}
    \item \textbf{Thiết kế 3D trong SolidWorks:} Xây dựng mô hình chi tiết gồm
    khung, bánh xe, giá đỡ, trục, động cơ và các liên kết cơ khí.

    \item \textbf{Đơn giản hóa mô hình:} Lược bỏ các chi tiết không ảnh hưởng
    đến chuyển động (lỗ bắt vít, góc vát,...) để giảm tải tính toán. Giữ nguyên
    các thông số quan trọng: kích thước tổng thể, vị trí khớp, bán kính bánh xe,
    khối lượng, trọng tâm và moment quán tính.

    \item \textbf{Xuất định dạng URDF:} URDF (Unified Robot Description Format)
    mô tả đầy đủ cấu trúc robot: các khâu, khớp nối, vị trí tương đối, khối
    lượng, moment quán tính, giới hạn chuyển động và hình dạng hiển thị.

    \item \textbf{Kiểm tra trong Isaac Sim:} Xác minh vị trí các khâu, hướng
    trục khớp, tâm quay bánh xe, thuộc tính va chạm và hệ số ma sát bánh xe --
    mặt sàn. Hai bánh xe được thiết lập là khớp xoay; các bộ phận còn lại là
    khớp cố định.

    \item \textbf{Huấn luyện trong Isaac Lab:} Chuyển mô hình sang Isaac Lab,
    định nghĩa môi trường RL (observation space, action space, reward function)
    và khởi động quá trình huấn luyện song song trên GPU.
\end{enumerate}
